\hypersetup{pageanchor=false}
\newgeometry{left=1.5cm,right=1.5cm,top=2cm,bottom=2cm}
\begin{titlepage}
\thispagestyle{empty}
\begin{tikzpicture}[remember picture,overlay,inner sep=0,outer sep=0]
     \draw[line width=4pt] ([xshift=-1.5cm,yshift=-2cm]current page.north east) coordinate (A)--([xshift=1.5cm,yshift=-2cm]current page.north west) coordinate(B)--([xshift=1.5cm,yshift=2cm]current page.south west) coordinate (C)--([xshift=-1.5cm,yshift=2cm]current page.south east) coordinate(D)--cycle;

     \draw ([yshift=0.5cm,xshift=-0.5cm]A)-- ([yshift=0.5cm,xshift=0.5cm]B)--
     ([yshift=-0.5cm,xshift=0.5cm]B) --([yshift=-0.5cm,xshift=-0.5cm]B)--([yshift=0.5cm,xshift=-0.5cm]C)--([yshift=0.5cm,xshift=0.5cm]C)--([yshift=-0.5cm,xshift=0.5cm]C)-- ([yshift=-0.5cm,xshift=-0.5cm]D)--([yshift=0.5cm,xshift=-0.5cm]D)--([yshift=0.5cm,xshift=0.5cm]D)--([yshift=-0.5cm,xshift=0.5cm]A)--([yshift=-0.5cm,xshift=-0.5cm]A)--([yshift=0.5cm,xshift=-0.5cm]A);


     \draw ([yshift=-0.3cm,xshift=0.3cm]A)-- ([yshift=-0.3cm,xshift=-0.3cm]B)--
     ([yshift=0.3cm,xshift=-0.3cm]B) --([yshift=0.3cm,xshift=0.3cm]B)--([yshift=-0.3cm,xshift=0.3cm]C)--([yshift=-0.3cm,xshift=-0.3cm]C)--([yshift=0.3cm,xshift=-0.3cm]C)-- ([yshift=0.3cm,xshift=0.3cm]D)--([yshift=-0.3cm,xshift=0.3cm]D)--([yshift=-0.3cm,xshift=-0.3cm]D)--([yshift=0.3cm,xshift=-0.3cm]A)--([yshift=0.3cm,xshift=0.3cm]A)--([yshift=-0.3cm,xshift=0.3cm]A);

\end{tikzpicture}\par
\centering
\vspace*{0.35cm}
ĐẠI HỌC QUỐC GIA THÀNH PHỐ HỒ CHÍ MINH\\
TRƯỜNG ĐẠI HỌC BÁCH KHOA\\
KHOA KHOA HỌC VÀ KỸ THUẬT MÁY TÍNH\par
\vspace{1cm}
\includegraphics[width=3cm]{figures/hcmut.png}\par
\vspace{1cm}
{\large\bfseries LUẬN VĂN THẠC SĨ\par}
\vspace{.7cm}
\begin{minipage}{.84\linewidth}
\centering\rule{\linewidth}{.5pt}\par
\vspace{.45cm}
{\LARGE\bfseries Bảo vệ tính riêng tư về quỹ đạo\\
cho người dùng dịch vụ dựa trên vị trí\par}
\vspace{.45cm}
\rule{\linewidth}{.5pt}
\end{minipage}\par
\vspace{.4cm}
\textbf{NGÀNH:} KHOA HỌC MÁY TÍNH\par
\vfill
\begin{tabular}{@{}ll@{}}
\textbf{Giảng viên hướng dẫn:} & TS. Phan Trọng Nhân\\
\textbf{Học viên thực hiện:} & Ngô Minh Hạnh\\
\textbf{Mã số học viên:} & 2470010
\end{tabular}\par
\vfill
{\footnotesize TP. HỒ CHÍ MINH, THÁNG 11/2026\par}
\vspace*{.9cm}
\end{titlepage}
\restoregeometry

\frontmatter
\hypersetup{pageanchor=true}
\prelimheading{Lời tựa}
Khi tìm một nhà thuốc trên bản đồ, tôi quan tâm đến việc có thể ghé nơi nào thuận đường.
Nhìn từ phía nhà cung cấp dịch vụ, cùng lượt tìm kiếm ấy lại cho biết một điều khác:
người dùng có thể đang ở đâu. Nếu các lần tìm kiếm được lưu lại theo thứ tự, liệu một
người chỉ nhìn thấy những vị trí đã làm mờ vẫn có thể đoán được đường đi và những nơi
người dùng thường lui tới? Câu hỏi ấy là điểm xuất phát của luận văn này.

Tôi tiếp cận bài toán từ phía người quan sát. Với mỗi tình huống, tôi xác định điều họ
muốn suy ra, phần hành trình họ được phép nhìn thấy và dữ liệu có thể dùng để kiểm tra
phỏng đoán. Cách đặt câu hỏi đó cũng giúp tôi nhìn lại phía người dùng: sau khi bảo vệ
quỹ đạo, điện thoại còn tìm được những địa điểm phù hợp với nhu cầu ban đầu hay không?

Tôi xin cảm ơn thầy Phan Trọng Nhân đã hướng dẫn tôi trong quá trình thực hiện luận văn.
Những trao đổi với thầy về cách so sánh các phương pháp và cách đánh giá chất lượng
dịch vụ đã giúp tôi xác định rõ hơn câu hỏi nghiên cứu và phạm vi của từng kết luận.

\prelimheading{Lời cam đoan}
Tôi chịu trách nhiệm về nội dung luận văn và cách trình bày các kết quả nghiên cứu
trong bản nộp. Các ý tưởng kế thừa, dữ liệu sử dụng và phương pháp đối chứng được
dẫn nguồn trong luận văn. Đối với các phương pháp được điều chỉnh để so sánh, tôi
trình bày rõ phần điều chỉnh và phạm vi diễn giải kết quả.
\par\vspace{1.5cm}
\hfill\begin{tabular}{c}Học viên\\[1.6cm]Ngô Minh Hạnh\end{tabular}

\prelimheading{Tóm tắt}
Khi một người dùng tìm địa điểm nhiều lần trên đường đi, nhà cung cấp có thể kết hợp tọa độ,
thời điểm và thứ tự các truy vấn để suy ra hành trình. Luận văn xem toàn bộ chuỗi trao đổi
trong một phiên là đối tượng cần bảo vệ. Mười kịch bản S1--S10 xác định điều bên quan sát
muốn suy ra và thông tin mà họ được phép sử dụng. Chất lượng dịch vụ được đánh giá theo các
địa điểm phù hợp mà người dùng thực sự nhận được tại vị trí thật.

Phương pháp GeoI-BQ v1 dùng cơ chế hàm mũ trên mạng đường để tạo điểm tham chiếu thỏa
geo-indistinguishability. Một phép kiểm tra có nhiễu quyết định khi nào cần cập nhật tham
chiếu. Từ lịch sử đã bảo vệ, thiết bị ước lượng phân phối vị trí để chọn một đoạn gồm ba
khung truy vấn. Các khung được gửi theo lịch công khai. Thiết bị hợp nhất các địa điểm
nhận về, lọc và xếp hạng theo nhu cầu riêng của người dùng.

Bốn mệnh đề liên kết lần lượt phép tạo tham chiếu, quyết định tái sử dụng, phân phối chọn
đoạn và việc công bố nhiều đoạn. Với các giả thiết về miền hỗ trợ chung, ngân sách, trạng
thái phiên và phản hồi của nhà cung cấp, ta thu được cận tỷ số xác suất cho toàn bộ chuỗi
quan sát. Cận này giới hạn mức thay đổi suy luận so với thông tin bên quan sát có từ trước.
Phân tích chất lượng dịch vụ chỉ ra điều kiện để mức tăng điểm lập kế hoạch chuyển thành
mức tăng chất lượng thật, cùng vai trò của bộ nhớ địa điểm tại thiết bị.

Thực nghiệm trên 48 nhóm tuyến mô phỏng so sánh GeoI-BQ với ba bản điều chỉnh từ các
nghiên cứu liên quan, GPS thật và phương án gửi trực tiếp tham chiếu REM. Trong tập kiểm
thử, năm tác vụ định vị chính có Hit100 từ 0 đến 0,188\%. Recall@5 trung bình toàn phiên
đạt \PrimaryRecall\%, với khoảng tin cậy 95\% \PrimaryCI\%, đổi lại 155 lượt truy vấn tọa
độ mỗi 600 giây. Danh mục tĩnh gần bão hòa ngay ở gói đầu và phương án nền công khai có
Recall nhỉnh hơn. Kết luận thực nghiệm vì vậy gắn với quyền truy cập dữ liệu, nhóm tuyến
và các mô hình suy luận đã đánh giá.

\noindent\textbf{Từ khóa:} riêng tư quỹ đạo, geo-indistinguishability, mạng đường,
truy vấn liên tục, địa điểm, hợp thành xác suất.

\clearpage
\begingroup\small\setstretch{1.05}\setlength{\parskip}{0pt}\tableofcontents\endgroup
\clearpage
\begingroup\small\setstretch{1.05}\setlength{\parskip}{0pt}\phantomsection\addcontentsline{toc}{chapter}{Danh mục hình}\listoffigures\endgroup
\clearpage
\begingroup\small\setstretch{1.05}\setlength{\parskip}{0pt}\phantomsection\addcontentsline{toc}{chapter}{Danh mục bảng}\listoftables\endgroup

\chapter*{Danh mục từ viết tắt}
\addcontentsline{toc}{chapter}{Danh mục từ viết tắt}
\begin{small}\setstretch{1.1}
\begin{longtable}{@{}P{2.1cm}P{13.4cm}@{}}
LBS & Location-based Services -- dịch vụ dựa trên vị trí.\\
LSP & Location Service Provider -- nhà cung cấp dịch vụ vị trí.\\
POI & Point of Interest -- địa điểm được quan tâm.\\
Geo-I & Geo-indistinguishability -- bảo đảm privacy theo khoảng cách vị trí.\\
REM & Road Exponential Mechanism -- kernel mũ trên miền đường công khai trong luận văn.\\
DP & Differential Privacy -- riêng tư vi sai.\\
SUMO & Simulation of Urban MObility -- mô phỏng giao thông vi mô.\\
FCD & Floating Car Data -- trạng thái phương tiện theo thời gian từ mô phỏng.\\
OSM & OpenStreetMap -- nguồn dữ liệu bản đồ mở.\\
GPS/GNSS & Hệ thống định vị vệ tinh; phân biệt GPS dùng bảo vệ và GPS dùng xếp hạng local.\\
PMF & Probability Mass Function -- hàm xác suất trên miền rời rạc.\\
LR & Likelihood Ratio -- tỷ số xác suất/khả năng của cùng output dưới hai input.\\
TV & Total Variation -- khoảng cách biến thiên toàn phần giữa hai phân phối.\\
CI & Confidence Interval -- khoảng tin cậy.\\
BA/AUC & Balanced Accuracy; diện tích dưới đường ROC.\\
MAE/EIE & Sai số tuyệt đối trung bình; sai số suy luận kỳ vọng theo contract tương ứng.\\
ASR/DER & Các metrics trong paper dummy-based; chỉ tính khi đáp ứng đúng định nghĩa nguồn.\\
kNN & k-Nearest Neighbors -- mô hình láng giềng gần nhất.\\
JSON & Định dạng bản tin ứng dụng được dùng trong phép đếm byte.\\
N/A & Không áp dụng hoặc chưa đủ đầu vào hợp lệ để tính; không mang giá trị 0.\\
\end{longtable}\end{small}

\chapter*{Danh mục ký hiệu}
\addcontentsline{toc}{chapter}{Danh mục ký hiệu}
\begin{small}\setstretch{1.08}
\begin{longtable}{@{}P{2.5cm}P{13cm}@{}}
\textbf{Ký hiệu} & \textbf{Ý nghĩa và đơn vị}\\\midrule
$X=(x_i)_i$ & Trajectory thật ở các mốc có thể đọc; tọa độ trong một hệ chiếu mét cố định.\\
$C$ & Ngữ cảnh công khai: bản đồ, catalogue, clock, cấu hình, provider law.\\
$V,a_v$ & Tập trạng thái đường công khai và tọa độ của trạng thái $v$.\\
$Z_t,E$ & Điểm tham chiếu đã bảo vệ; lịch sử nội bộ gồm các Z và nhánh test.\\
$b_t(s)$ & Belief gần đúng về trạng thái người dùng, tính từ thông tin đã bảo vệ.\\
$Q_t,K$ & Khung gồm $K$ query slots gửi tại một mốc; prototype dùng $K=5$.\\
$A_a,\mathcal A_j$ & Một đoạn ba khung Q và thư viện các đoạn ở lần chọn $j$.\\
$\mathcal H_j$ & Lịch sử phân tích có đoạn đã chọn, lane giả và replies; mạnh hơn wire tọa độ.\\
$O_n$ & Transcript server sau $n$ lần chọn đoạn: Q, thứ tự, clock, request/reply, retry.\\
$w_a,p_a$ & Phân phối nền và phân phối thích nghi khi chọn đoạn.\\
$v_a,z_a$ & Điểm lập kế hoạch trước chuẩn hóa; điểm bị chặn đưa vào kernel.\\
$u,C_s$ & Chi phí một đơn vị và cap Geo-I từng phiên, đơn vị m$^{-1}$.\\
$B,H_{\rm cfg},M$ & Tham số ngân sách danh nghĩa, tham số accounting và số đơn vị cap; $M=2H_{\rm cfg}-1$.\\
$\varepsilon_{Q,j},\Gamma_Q$ & Allowance chọn đoạn thứ $j$ và cap Q cả phiên, không có đơn vị.\\
$L,k$ & Độ sâu phản hồi mỗi loại POI/tọa độ; độ sâu reference đánh giá, lần lượt 60 và 5.\\
$\psi_t,U_t,T_{t,\psi}$ & Nhu cầu thật tại local; cache hợp lệ; tập POI reference tại GPS/nhu cầu đang xét.\\
$D_\infty(X,X')$ & Độ lệch tọa độ lớn nhất giữa hai traces trên cùng lịch, đơn vị mét.\\
$R_n,\delta_n$ & Cận LR của kênh Q và cận TV cả phiên so với matched public base.\\
$\eta,\tau,\nu$ & Sai số calibration/proxy và chuyển miền có điều kiện; chưa biết không được đặt bằng 0.\\
\end{longtable}\end{small}
